\documentclass{article}
\usepackage[margin=.75in]{geometry}
\usepackage{bookmark}
\usepackage{graphicx} % Required for inserting images
\usepackage{float}
\usepackage[titletoc,toc,page]{appendix}
\usepackage{wrapfig}
\usepackage{caption}
\usepackage{subcaption}
\usepackage{mathtools}
\usepackage{authblk}
\usepackage{hyperref}

\hypersetup{
    colorlinks,
    citecolor=blue,
    filecolor=blue,
    linkcolor=blue,
    urlcolor=blue
}

\floatstyle{plain} 
\restylefloat{figure}
\title{AST326 Lab 3: Type Ia Supernova}
\author{Luke Wriglesworth \\ \small luke.wriglesworth@mail.utoronto.ca}
\date{\today}
\begin{document} % Add this line
\maketitle

\begin{abstract}
    In this report, aperture and differential photometry techniques are used to analyze a type Ia supernova, AST325-326-SN. Photometric calibration is 
    done using three well studied reference stars of known magnitudes and uncertainties, and differential photometry is used to plot the magnitude of 
    AST325-326-SN over measurements ranging from September 20, 2015, to October 30, 2015. Statistical methods are used to filter out poor quality images, and 
    the peak magnitude of AST325-326-SN is found to be $\approx 18.68 \pm 0.36$ on 2015-10-30, after its initial detection on 2015-9-24 with a signal-noise ratio $>$ 3. By 
    applying power law fitting to the early stages of the normalized light curve of AST325-326-SN, the epoch of first light is estimated to be $9.59 \pm 5.997$ 
    days prior to the epoch of first detection. 
\end{abstract}


\section{Introduction}
Aperture and differential photometry are two very important tools for astronomers, allowing accurate measurements to be made of point light sources. Aperture photometry is a technique used to measure the intensity of a point light source
by summing the intensities of all the pixels inside a circle of a given radius over the source, and subtracting a measurement of noise taken of the surrounding sky background. Differential photometry is a related 
technique used to convert these intensity measurements into magnitudes. This is done by using the known magnitude of a reference light source
in the frame, and comparing it to the measurement of the target light source. These techniques can be used on just a single image, but they are much more effective over a range of measurements
taken throughout long time periods. In this lab, these techniques are used to analyze a type Ia supernova, AST325-326-SN, and determine its epoch of first light, epoch of first detection, epoch of 
last detection, and its peak magnitude.
    
A set of 433 images of AST325-326-SN are analyzed, taken from September 20, 2015, to October 30, 2015. The images are taken in the B-band range and are of varying levels of quality. Before
analysis is done, the images are first filtered based on the median value of all the images, and then they are filtered based on intensity measurements of reference stars and measurements of 
AST325-326-SN. The remaining images are then analyzed using aperture and differential photometry techniques, and the results are used to plot the magnitude of AST325-326-SN over time. To better
estimate the peak magnitude of AST325-326-SN, polynomial fitting is applied to the weighted mean of the distribution. Lastly, power law fitting is applied to the early stages of the normalized 
light curve of AST325-326-SN to estimate the epoch of first light.

All the work in this lab was done independently, and all the code and plots were generated by the individual lab members. We only shared our results with each other throughout the process to make sure everyone 
was on the right track.
\newpage

\section{Data and Observations}

The data used in this report is obtained between September 20, 2015, and October 30, 2015. The measurements are in the form of 433, $710\times850$ pixel CCD .fits files in the B-band
range, of varying degrees of quality. Basic noise reduction techniques such as bias and dark subtraction, and flat fielding were already applied to the images prior to analysis. 
Initially, AST325-326-SN is not visible in the images, and gradually becomes visually apparent over the course of the measurement period. In every image that is of reasonable quality, 
three reference stars are visible, and can be used for photometrical calibration and to perform differential photometry on AST325-326-SN. The reference stars are labelled as REF 1, 
REF 2 and REF 3. Their sky coordinates, as well as their B-band magnitudes and uncertainties are displayed in Table \ref{tab:ref_stars}.

\begin{table}[htbp]
    \centering
    \begin{tabular}{c|c|c}
        \hline
        \textbf{Reference Star} & \textbf{(RA, DEC)} & \textbf{B-Band Magnitudes} \\
        \hline
        REF 1 & (00:56:49.70, -37:01:38.31) & $16.32 \pm 0.07$ \\
        REF 2 & (00:56:46.43, -37:02:29.50) & $17.16 \pm 0.08$  \\
        REF 3 & (00:56:58.27, -36:58:16.60) & $15.61 \pm 0.02$ \\
        \hline
    \end{tabular}
    \caption{Reference Stars Used for Photometric Calibration and Differential Photometry}
    \label{tab:ref_stars}
\end{table}

\section{Data Reduction and Methods}
\subsection{Filtering Poor Quality Images}

The raw data from the observations are of a very wide range in quality. Many of the images are dominated by technical glitches and noise, and some are
of good quality, but completely off center from the targets. To account for this, three different filtering techniques are applied to the images at
different points in the data reduction process. One major initial issue is overflowing pixels. Some very bright points in the images exceed the maximum
intensity value that the CCD can record, and cause the CCD to compensate by recording these pixels as negative values. To account for this, numerous images were 
plotted as 1D histograms, and a general threshold for negative pixels is determined that still permits important sub-zero noise values. Each image with pixels below 
this threshold will have the pixels wrapped back around to positive values. Next, the median value of each pixel in each image was calculated, and the residuals of each image in 
relation to this median value were plotted on a scatter plot. This allows an arbitrary boundary range to be set to filter out any images differing wildly from the median image.

\begin{figure}[htbp]
    \centering
    \includegraphics[width=0.8\textwidth]{/home/luke/Documents/Course Content/University of Toronto/Year 3/AST326/AST326-labs/Lab 3/plt/median.pdf}
    \caption{Plot of each image's distance from the median image.}
    \label{fig:median_filter}
\end{figure}

This median filtering technique works well at removing images that are obviously low quality, but there are still some sneakier images where either the reference stars are not visible, 
or the measurements of AST325-326-SN return unreasonably high or low intensities. To account for this, two more filters were put in place during the aperture photometry measurement
process. Prior to recording any measurement, the intensities of REF1 and REF2 are compared to an arbitrary threshold value, determined by the intensity values of the majority of the measurements.
Any image that does not meet this reference star intensity threshold gets discarded. Finally, the intensity of AST325-326-SN is checked as a final filter, and if it returns \texttt{nan} or \texttt{inf}, 
the image will be discarded, even if the reference stars had good measurements. At the end of the whole filtering process, the 433 images are reduced to 313 images.

\subsection{Aperture Photometry}

\begin{figure}[H]
    \centering
    \begin{subfigure}[b]{0.5\textwidth}
        \includegraphics[width=\textwidth]{/home/luke/Documents/Course Content/University of Toronto/Year 3/AST326/AST326-labs/Lab 3/plt/centroid.pdf}
        \caption{An example of a centroid calculation applied to REF 1.}
        \label{fig:centroid}
    \end{subfigure}
    \hfill
    \begin{subfigure}[b]{0.65\textwidth}
        \includegraphics[width=\textwidth]{/home/luke/Documents/Course Content/University of Toronto/Year 3/AST326/AST326-labs/Lab 3/plt/gaussian_fit.pdf}
        \caption{A Gaussian function being fitted to the central row of the point spread function of REF 1.}
        \label{fig:gaussian}
    \end{subfigure}
    \caption{Centroid calculation and Gaussian fit}
\end{figure}

For measurements that pass all the filtering criteria, measurements are recorded using a typical aperture photometry pipeline. Each image's header file is read in a loop, and the sky coordinates of the reference
stars are extracted and converted to pixel coordinates. Next, centroids are calculated for each of the reference stars (Figure \ref{fig:centroid}), and this centroid coordinate is used to fit a Gaussian function to a 1D cutout of the 
stars point spread function (Figure \ref{fig:gaussian}). The radius of the aperture is set to be 3 times the standard deviation of the Gaussian fit, and the sky annulus radii are arbitrarily set using a formula common to all the stars, 
using the aperture radius as a parameter.

The aperture and annulus functions are both implemented in similar ways, using a binary mask to select all pixels within a specific radius or radii. The measurements in the aperture are taken to be 
the summation of all pixels within the aperture mask, and the measurements in the sky annulus are taken to be the average value per pixel, multiplied by the area of the annulus. An example of these functions in action is shown in 
Figures \ref{fig:aperture} and \ref{fig:annulus}. The uncertainty in the aperture measurements 
is taken to be the square root of each pixel intensity divided by the gain of the instrument, which is taken to be 4 in this lab. The uncertainty in the sky annulus measurements is taken to be the square root 
of the number of pixels in the annulus, multiplied by the standard deviation of the pixels in the sky annulus. The final intensity measurement of each image is taken to be the difference between the aperture intensity measurement 
and the background noise measurement captured by the sky annulus. 

\begin{figure}[H]
    \centering
    \begin{subfigure}[b]{0.45\textwidth}
        \includegraphics[width=\textwidth]{/home/luke/Documents/Course Content/University of Toronto/Year 3/AST326/AST326-labs/Lab 3/plt/aperture.pdf}
        \caption{An aperture measurement being done on REF 1.}
        \label{fig:aperture}
    \end{subfigure}
    \hfill
    \begin{subfigure}[b]{0.45\textwidth}
        \includegraphics[width=\textwidth]{/home/luke/Documents/Course Content/University of Toronto/Year 3/AST326/AST326-labs/Lab 3/plt/annulus.pdf}
        \caption{An annulus measurement being done on REF 1.}
        \label{fig:annulus}
    \end{subfigure}
    \caption{Aperture and sky annulus measurements being done on REF 1. Here, the lighter portions of each image are being measured and darker parts excluded.}
\end{figure}

With the intensity signal from the aperture and the noise signal from the sky annulus, the signal-noise ratio of each measurement can be calculated, which will allow 
a statistical threshold to be used for the epoch of first detection of AST325-326-SN, which is taken to be 3 based on recommendations from Professor Moon. Based on the lecture notes, the 
propagated uncertainty from the aperture and annulus for a function $f(x,y) = \frac{x}{y}$ is given by Formula \ref{math:uncertainty}, where $f$ is the signal-noise ratio, 
$x$ are the aperture measurements, and $y$ are the sky annulus measurements:

\begin{gather}
    \label{math:uncertainty}
        (\frac{\sigma_f}{f})^2 = (\frac{\sigma_x}{x})^2 + (\frac{\sigma_y}{y})^2
\end{gather}

The signal-noise ratio was found to first exceed three in the image \textbf{2015-09-24 21:33:29}, with a signal-noise ratio uncertainty of $\approx 1.65$. This is taken to be the epoch of first detection for AST325-326-SN. 
Additionally, the signal noise ratio was found to drop back below three in the image \textbf{2015-10-30 01:34:52}, with a signal-noise ratio uncertainty of $\approx 0.101$. This is taken to be the epoch of last detection for AST325-326-SN.

\subsection{Differential Photometry}

Next, the recorded intensities of AST325-326-SN are converted into magnitudes using differential photometry. Before making any calculations with AST325-326-SN, the magnitudes of the reference stars
are calculated and compared with their expected magnitudes provided in the lab document. The magnitude residuals of REF 1 and REF 2 are plotted in Figure \ref{fig:ref_residuals} compared with their expected magnitudes. 
Ref 3 is not included, as it is used as a reference to calculate the magnitudes of REF 1 and REF 2. As you can see from Figure \ref{fig:ref_residuals}, the residuals are quite small, so good results are expected from the setup. 

\begin{figure}[H]
    \centering
    \includegraphics[width=0.4\textwidth]{/home/luke/Documents/Course Content/University of Toronto/Year 3/AST326/AST326-labs/Lab 3/plt/residuals.pdf}
    \caption{Residuals of the reference stars compared to their expected magnitudes.}
    \label{fig:ref_residuals}
\end{figure}

Next, the magnitudes of AST325-326-SN can be plotted over time, using the magnitude of REF 3 as a reference: 

\begin{gather}
    m_{SN} = 2.5\log_{10}(\frac{I_3}{I_{SN}}) + m_{3}
\end{gather}

Here, $I_3$ is the intensity of REF 3, $I_{SN}$ is the intensity of AST325-326-SN, and $m_3$ is the magnitude of REF 3. The magnitude of REF 3 is taken to be $15.61 \pm 0.02$, as given in Table \ref{tab:ref_stars}. 
The uncertainties associated with magnitudes of AST325-326-SN can be propagated at this point from the three reference stars. For example, the uncertainties associated with REF 1, REF2, and REF 3, are: 

\begin{gather}
    \sigma_{1} = (\frac{2.5}{\ln(10)})(\frac{\sigma_{I_{SN}}}{I_{SN}}) \\
    \sigma_{2} = (\frac{2.5}{\ln(10)})(\frac{\sigma_{I_{2}}}{I_{2}}) \\
    \sigma_{3} = 0.02
\end{gather}

So the total propagated error through the magnitude calculations are:

\begin{gather}
    \label{math:propogated_error}
    \sigma_{SN} = \sqrt{\sigma_{1}^2 + \sigma_{2}^2 + \sigma_{3}^2}
\end{gather}

A plot of the magnitudes of AST325-326-SN and the propagated uncertainty are shown in Figure \ref{fig:raw_plot}.

\begin{figure}[H]
    \centering
    \includegraphics[width=.85\textwidth]{/home/luke/Documents/Course Content/University of Toronto/Year 3/AST326/AST326-labs/Lab 3/plt/raw_mag.pdf}
    \caption{Magnitudes of AST325-326-SN as a function of image number.}
    \label{fig:raw_plot}
\end{figure}
\section{Data Analysis and Modeling}

As it is difficult to further analyze this plot as a function of image number, it is binned per day and shown in Figure \ref{fig:binned_plot}. 

\begin{figure}[H]
    \centering
    \includegraphics[width=.80\textwidth]{/home/luke/Documents/Course Content/University of Toronto/Year 3/AST326/AST326-labs/Lab 3/plt/binned_uncertainty.pdf}
    \caption{Magnitudes of AST325-326-SN as a function of date.}
    \label{fig:binned_plot}
\end{figure}

There is still a lot of noise going on this plot, and since the goal is to find the peak magnitude of AST325-326-SN, it is useful to find the weighted average of the distribution, so that the outlying measurements are given
less weight. In Figure \ref{fig:weighted_average}, the weighted average of the distribution is shown, where the weights are the inverse of the uncertainty in each measurement. 

\begin{figure}[H]
    \centering
    \includegraphics[width=.90\textwidth]{/home/luke/Documents/Course Content/University of Toronto/Year 3/AST326/AST326-labs/Lab 3/plt/weighted_mean.pdf}
    \caption{Weighted average of the magnitudes of AST325-326-SN as a function of date.}
    \label{fig:weighted_average}
\end{figure}
\subsection{Polynomial Fitting}

To find the peak magnitude of AST325-326-SN, a polynomial fitting is applied to the weighted average of the magnitudes as a function of date. Fittings for 
polynomials of degrees 1 through 4 were applied, and the degree 4 polynomial was found to provide the lowest uncertainty in the peak magnitude calculation without overfitting
the distribution. Using the weighted mean fourth degree polynomial, AST325-326-SN is found to have a peak magnitude of $18.68 \pm 0.37$. A plot of the polynomial fitting is shown 
in Figure \ref{fig:poly_fit}.

\begin{figure}[H]
    \centering
    \includegraphics[width=1\textwidth]{/home/luke/Documents/Course Content/University of Toronto/Year 3/AST326/AST326-labs/Lab 3/plt/weighted_mean_poly.pdf}
    \caption{Fourth degree polynomial fitting of AS.}
    \label{fig:poly_fit}
\end{figure}

\subsection{Power Law Fitting}
Lastly, a power law fitting is applied to the early stages of the normalized light curve of AST325-326-SN to attempt to estimate the epoch of first light.
A threshold of 40 percent of the peak normalized intensity is used as the upper limit of the fitting. The normalized intensity values of AST325-326-SN are plotted
in Figure \ref{fig:power_fit}, with a function $f(t)=C(t-t_1)^2$ fitted to it. The fitting returned optimal values of $C=0.00275 \pm 1.45\times10^{-4}$, $t=-9.59 \pm 5.997$, 
indicating a poor estimate of the epoch of first light. The fit also gave a reduced $\chi^2$ value of $\approx 0.0053$ which despite being very low, is likely an indication 
of under-fitting of the data. Despite this setback, an estimate of the epoch of first light can still be made, which would be somewhere between 4 and 15 days prior to the epoch of first detection.

\begin{figure}[H]
    \centering
    \includegraphics[width=.80\textwidth]{/home/luke/Documents/Course Content/University of Toronto/Year 3/AST326/AST326-labs/Lab 3/plt/power_law.pdf}
    \caption{Power law fitting applied to the early stages of the normalized light curve of AST325-326-SN.}
    \label{fig:power_fit}
\end{figure}

\section{Discussion}

By applying aperture and differential photometry techniques, a great amount of information about AST325-326-SN has been determined. 
These results are summarized in Table \ref{tab:results}. 

\begin{table}[H]
    \centering
    \begin{tabular}{|c|c|}
        \hline
        \textbf{Result} & \textbf{Value} \\
        \hline
        Epoch of First Detection & 2015-09-24 21:33:29 ($\pm$ 1.65 s/n ratio)\\
        \hline
        Epoch of Peak Brightness &  2015-10-08 06:14:07 ($\pm$ 0.37 magnitude)\\
        \hline
        Epoch of First Light &  $-9.59 \pm 5.997$ days\\
        \hline
        Peak Magnitude & $18.68 \pm 0.37$ \\
        \hline
        Epoch of Non-Detection & 2015-10-30 01:34:52 ($\pm$ 0.101 s/n ratio)\\
        \hline
    \end{tabular}
    \caption{Summary of Results}
    \label{tab:results}
\end{table}

The magnitude result seem to be very reasonable, as it is in line with typical type Ia supernovae which are known to have peak magnitudes of $\approx 19.5$. The epoch of first light result however, is most likely not accurate. I had 
great difficulties getting a decent power law fitting, and it may be due to not having enough data points in the early stages of the light curve. The image filtering process may have been too aggressive, and removed important images that are needed
to complete the fitting. Nevertheless, from visual inspection of the curve, I believe the epoch of first light is still within the uncertainty range of the calculation, and is most likely somewhere between 4 and 9 days prior to the epoch of first detection.

\newpage

\appendix

\begin{appendices}
\section{Python Code}
\subsection{Image Filtering}
\begin{verbatim}
    def median_filter_images(images, names, distances, upper, lower):
    good = []
    good_i = []
    bad = []
    bad_i = []
    for i, image in enumerate(images):
        if not ((distances[i] >= upper) or (distances[i] <= lower)):
            good.append(image)
            good_i.append(i)
        else:
            bad.append(image)
            bad_i.append(i)
    print("{g} good images out of {t}".format(g=len(good),t=len(images)))
    return np.array(good), np.array(good_i), np.array(bad), np.array(bad_i)   
\end{verbatim}

\subsection{Reading / Sorting / Viewing Images, Fixing Overflowed Pixels}

\begin{verbatim}
    def read_file(filename):
    # Need to fix bright pixels wrapping around into negative values.
    image = fits.open("AST325-326-SN-short/{f}".format(f=filename))[0].data.astype(int)
    overflow_magnitude = np.abs(image[image < -500])
    max_value = 2**15-1
    
    # unwrap negative pixels to positive
    image[image < -500] = (max_value-overflow_magnitude) + max_value
    return image

def get_header(filename):
    header = fits.open("AST325-326-SN-short/{f}".format(f=filename))[0].header
    return header

def get_headers(namelist):
    headers = []
    for i, name in enumerate(namelist):
        current_header = get_header(name)
        headers.append(current_header)
    return headers

# Puts the image names in order and returns a new ordered namelist
def sort_names(names):
    name_list = np.loadtxt((names), dtype=str)
    return np.sort(name_list)

# Create a numpy array of images, based on a name list
def image_array(namelist):
    images = []
    for i, name in enumerate(namelist):
        images.append(read_file(name))
    return np.array(images)
    
# View all images in array, also can plot histograms to check for overflow
def view_all(images, names,xcen,ycen,box,zoom=False):
    headers = get_headers(names)
    for i, image in enumerate(images):
        w = WCS(headers[i])
        fig = plt.figure(figsize=(10,10))
        ax = fig.add_subplot(111, projection=w)
        print(names[i])
        print(headers[i]['DATE-OBS'])
        #print("Max Value: {max}, Min Value: {min}".format(max=np.max(image),min=np.min(image)))
        if zoom:
            ax.set_xlim(xcen-box,xcen+box)
            ax.set_ylim(ycen-box,ycen+box)
        ax.imshow(image, vmin=-100, vmax = 1e2, origin = 'lower', cmap = 'inferno')
        plt.show()
        
        # # Plot histogram of entire image as one row to check for overflowing pixels
        # plt.figure(figsize=(12,5))
        # plt.hist(image.flatten(), bins = 'sturges', color = 'k', histtype='step')
        # plt.grid(alpha=0.5)
        # plt.yscale('log')
        # plt.xlabel("Pixel Value [ADU]", size = 16)
        # plt.ylabel("Number of Pixels", size = 16)
        # plt.show()
        
# View a single image
def view(image):
    plt.figure(figsize=(10,10))
    plt.imshow(image, vmin=-100, vmax = 1e2, origin = 'lower', cmap = 'inferno')
    plt.show()

\end{verbatim}
\subsection{Aperture and Annulus}
\begin{verbatim}

    def aperture(image, xcenter, ycenter, radius, plot=False):
    ysize, xsize = image.shape
    y, x = np.mgrid[:ysize, :xsize]
    mask = ((x - xcenter)**2 + (y - ycenter)**2) <= radius**2
    aperture_sum = np.sum(image * mask)
    if plot:
        plt.xlim(xcenter-20,xcenter+20)
        plt.ylim(ycenter-20,ycenter+20)
        plt.imshow(image, vmin=-100, vmax = 1e2, origin='lower', cmap='inferno')
        plt.imshow(mask, origin='lower', cmap='inferno', alpha=0.6)
        plt.savefig("plt/aperture.pdf")
        plt.show()
    return aperture_sum

def annulus(image, xcenter, ycenter, inner_radius, outer_radius, plot=False):
    ysize, xsize = image.shape
    y, x = np.mgrid[:ysize, :xsize]
    r_squared = (x - xcenter)**2 + (y - ycenter)**2
    mask = (r_squared <= outer_radius**2) & (r_squared >= inner_radius**2)
    annulus_sum = np.sum(image * mask)
    num_pixels = np.sum(mask)
    level_per_pixel = annulus_sum / num_pixels if num_pixels else 0
    result = level_per_pixel * num_pixels
    
    # Uncertainty
    sigma_bg = np.std(image*mask)
    annulus_uncertainty = np.sqrt(num_pixels * sigma_bg)
    
    if plot:
        plt.xlim(xcenter-20,xcenter+20)
        plt.ylim(ycenter-20,ycenter+20)
        plt.imshow(image, vmin=-100, vmax = 1e2, origin='lower', cmap='inferno')
        plt.imshow(mask, origin='lower', cmap='gray', alpha=0.6)
        plt.savefig("plt/annulus.pdf")
        plt.show()
    return result, annulus_uncertainty
\end{verbatim}

\subsection{Finding Supernova Aperture Radius}
\begin{verbatim}
    # Given a range of sample images, returns the radius to use for the aperture & annulus, based on 3
    # times the value of the std of the fitted gausian to each image. 
    def sn_radius(sample_range1,sample_range2,x,y,box_sn):
        gauss_std1 = []
        gauss_std2 = []
        for i,image in enumerate(sample_range1):
            sn_x_cen,sn_y_cen = centroid(x,y,image,box_sn)
            cutout_sn = image[int(sn_y_cen-box):int(sn_y_cen+box), int(sn_x_cen-box):int(sn_x_cen+box)]
            x_sn_gauss, y_sn_gauss = np.arange(cutout_sn.shape[0]), cutout_sn[cutout_sn.shape[0]//2, :]
            popt_sn, pcov_sn = curve_fit(gaus, x_sn_gauss, y_sn_gauss, p0=[np.max(cutout_sn), cutout_sn.shape[0]//2, 10, 0])
            gauss_std1.append(abs(popt_sn[2]))
        for i,image in enumerate(sample_range2):
            sn_x_cen,sn_y_cen = centroid(x,y,image,box_sn)
            cutout_sn = image[int(sn_y_cen-box):int(sn_y_cen+box), int(sn_x_cen-box):int(sn_x_cen+box)]
            x_sn_gauss, y_sn_gauss = np.arange(cutout_sn.shape[0]), cutout_sn[cutout_sn.shape[0]//2, :]
            popt_sn, pcov_sn = curve_fit(gaus, x_sn_gauss, y_sn_gauss, p0=[np.max(cutout_sn), cutout_sn.shape[0]//2, 10, 0])
            gauss_std2.append(abs(popt_sn[2]))
        avg_std1 = np.mean(gauss_std1)
        avg_std2 = np.mean(gauss_std2)
        radius_sn = 3*np.mean((avg_std1+avg_std2)/2)
        sn_r1 = radius_sn + radius_sn / 3
        sn_r2 = sn_r1 + sn_r1 / 2
        print(radius_sn)
        return radius_sn, sn_r1, sn_r2
\end{verbatim}

\subsection{Main Program Loop}
\begin{verbatim}
    # Sort image names in order
    sorted_names = sort_names("AST325-326-SN-filelist.txt")
    
    # Get numpy array of sorted images
    images = image_array(sorted_names)
    
    # View images in order 
    #view_all(images, sorted_names)
    
    # Need to remove bad images...
    median = np.median(images)
    image_residuals = image_residuals(images,median)
    
    # Finding the sum-squared distance from the median of the images
    distances = sum_residuals(image_residuals)
    median_error = np.median(distances)
    mad = sp.stats.median_abs_deviation(distances)
    
    upper_bound = 1.3e11
    lower_bound = 5.2e9
    
    # # Plot the distances to find outliers and set a threshold for images to keep. 
    image_numbers = np.arange(0,len(sorted_names))
    plt.title("Distances From Median Image")
    plt.xlabel("Image Number")
    plt.ylabel("Distance")
    plt.yscale("log")
    plt.axhline(median_error, color = "Red", ls="--", label="Median Image")
    plt.axhline(upper_bound, color="green", label="Upper Threshold: {e}".format(e=upper_bound))
    plt.axhline(lower_bound, color="brown", label="Lower Threshold: {e}".format(e=lower_bound))
    plt.scatter(image_numbers,distances,label="Distance", s=0.5)
    plt.grid(".-")
    plt.legend(loc="best")
    plt.savefig("plt/median.pdf")
    plt.show()
    
    # Find bad/good images 
    good_imgs, good_indices, bad_imgs, bad_indices = median_filter_images(images, sorted_names, distances, upper_bound, lower_bound)
    good_names = sorted_names[good_indices]
    
    #view_all(good_imgs,good_names,0,0,0)
    #view_all(bad_imgs,sorted_names[bad_indices])
    #aperture(good_imgs[0],250,250,20,plot=True)
    
    # Reference Star 1: (00:56:49.70, -37:01:38.31)
    # Reference Star 2: (00:56:46.43, -37:02:29.50)
    # Reference Star 3: (00:56:58.27,-36:58:16.60)
    
    ref_1 = SkyCoord('00:56:49.70 -37:01:38.31', unit=('hourangle', 'deg'))
    ref_2 = SkyCoord('00:56:46.43 -37:02:29.50', unit=('hourangle', 'deg'))
    ref_3 = SkyCoord('00:56:58.27 -36:58:16.60', unit=('hourangle', 'deg'))
    sn = SkyCoord('00:57:03.19 -37:02:23.6', unit=('hourangle', 'deg'))
    
    box=15
    box_sn=10
    measurement_headers = []
    aperture_measurement_ref1 = []
    aperture_measurement_ref2 = []
    aperture_measurement_ref3 = []
    aperture_measurement_sn = []
    annulus_measurement_ref1 = []
    annulus_measurement_ref2 = []
    annulus_measurement_ref3 = []
    annulus_measurement_sn = []
    
    # Annulus uncertainties
    annulus_uncertainty_ref1 = []
    annulus_uncertainty_ref2 = []  
    annulus_uncertainty_ref3 = []
    annulus_uncertainty_sn = []
    
    
    measured_indices_in_good_imgs = []
    
    sample1 = good_imgs[150:160]
    sample2 = good_imgs[250:260]
    
    #view_all(sample1,good_names[150:160],100,125,box_sn,zoom=True)
    #print("Second sample...")
    #view_all(sample2,good_names[250:260],100,125,box_sn,zoom=True)
    
    radius_sn, sn_r1, sn_r2 = sn_radius(sample1,sample2,100,125,box_sn)
    
    for i,image in enumerate(good_imgs):
        header = get_header(good_names[i])
        w = WCS(header)
        # Pixel positions of objects
        x1, y1 = w.all_world2pix(ref_1.ra.deg, ref_1.dec.deg, 1)
        x2, y2 = w.all_world2pix(ref_2.ra.deg, ref_2.dec.deg, 1)
        x3, y3 = w.all_world2pix(ref_3.ra.deg, ref_3.dec.deg, 1)
        #x_sn,y_sn = w.all_world2pix(sn.ra.deg, sn.dec.deg, 1)
        x_sn,y_sn = 100,125
    
        # Centroids of objects
        x1_cen,y1_cen = centroid(x1,y1,image,box)
        x2_cen,y2_cen = centroid(x2,y2,image,box)
        x3_cen,y3_cen = centroid(x3,y3,image,box)
        sn_x_cen,sn_y_cen = centroid(x_sn,y_sn,image,box_sn)
        
        # I'm checking the intensity of reference star 2 and 1 to skip any images with unreasonably low values, since
        # some images with missing stars made it past the initial median filter.     
        initial_measurement_ref2 = aperture(image,x2_cen,y2_cen,5)
        initial_measurement_ref1 = aperture(image,x1_cen,y1_cen,5)
        initial_measurement_ref3 = aperture(image,x3_cen,y3_cen,5)
        initial_measurement_sn = aperture(image,sn_x_cen,sn_y_cen,5)
        threshold = 2000
        if initial_measurement_ref2 > threshold and initial_measurement_ref1 > threshold and initial_measurement_ref3 > threshold:        
            # Create a cutout around the centroid positions
            cutout_ref1 = image[int(y1_cen-box):int(y1_cen+box), int(x1_cen-box):int(x1_cen+box)]
            cutout_ref2 = image[int(y2_cen-box):int(y2_cen+box), int(x2_cen-box):int(x2_cen+box)]
            cutout_ref3 = image[int(y3_cen-box):int(y3_cen+box), int(x3_cen-box):int(x3_cen+box)]
    
            # Fit a Gaussian to the central row of the cutout to find aperture radius
            x_ref1, y_ref1 = np.arange(cutout_ref1.shape[0]), cutout_ref1[cutout_ref1.shape[0]//2, :]
            x_ref2, y_ref2 = np.arange(cutout_ref2.shape[0]), cutout_ref1[cutout_ref2.shape[0]//2, :]
            x_ref3, y_ref3 = np.arange(cutout_ref3.shape[0]), cutout_ref1[cutout_ref3.shape[0]//2, :]
    
            popt_ref1, pcov_ref1 = curve_fit(gaus, x_ref1, y_ref1, p0=[np.max(cutout_ref1), cutout_ref1.shape[0]//2, 10, 0])
            popt_ref2, pcov_ref2 = curve_fit(gaus, x_ref2, y_ref2, p0=[np.max(cutout_ref2), cutout_ref2.shape[0]//2, 10, 0])
            popt_ref3, pcov_ref3 = curve_fit(gaus, x_ref3, y_ref3, p0=[np.max(cutout_ref3), cutout_ref3.shape[0]//2, 10, 0])
    
            # # Plot the central row and the fitted Gaussian
            # plt.figure(figsize=(10, 6))
            # plt.plot(x_ref1, y_ref1, 'x', color='k', label='Central Row of REF 1')
            # plt.plot(x_ref1, gaus(x_ref1, *popt_ref1), color='r', alpha=0.3, label=f'Best-fit Gaussian (STD {round(popt_ref1[2], 1)})')
            # plt.legend(loc="best", fontsize=14)
            # plt.grid(".-")
            # plt.savefig("plt/gaussian_fit.pdf")
            # plt.show()
            
            # lets use 3 times the STD of the gaussian for the aperture radius. 
            radius_ref1 = abs(3*popt_ref1[2])
            radius_ref2 = abs(3*popt_ref2[2])
            radius_ref3 = abs(3*popt_ref3[2])
    
            # annulus radii
            ref1_r1 = radius_ref1 + radius_ref1 / 3
            ref1_r2 = ref1_r1 + ref1_r1 / 2
            ref2_r1 = radius_ref2 + radius_ref2 / 3
            ref2_r2 = ref2_r1 + ref2_r1 / 2
            ref3_r1 = radius_ref3 + radius_ref3 / 3
            ref3_r2 = ref3_r1 + ref3_r1 / 2
    
            aperture1 = aperture(image, x1_cen, y1_cen, radius_ref1)
            aperture2 = aperture(image, x2_cen, y2_cen, radius_ref2)
            aperture3 = aperture(image, x3_cen, y3_cen, radius_ref3)
            aperture_sn = aperture(image, sn_x_cen, sn_y_cen, radius_sn)
            
            # Check again that we are measuring valid values, including supernova data
            if not np.isnan(aperture1) and not np.isnan(aperture2) \
            and not np.isnan(aperture3) and not np.isnan(aperture_sn):
    
                if aperture1 > threshold and aperture2 > threshold \
                and aperture3 > threshold and aperture_sn > 0:
                            
                    # Record index
                    measured_indices_in_good_imgs.append(i)
                    
                    # Record measurements
                    aperture_measurement_ref1.append(aperture(image, x1_cen,y1_cen,radius_ref1, plot=False))
                    annulus_measurement_ref1.append(annulus(image,x1_cen,y1_cen,ref1_r1,ref1_r2, plot=False))
                    aperture_measurement_ref2.append(aperture(image, x2_cen,y2_cen,radius_ref2))
                    annulus_measurement_ref2.append(annulus(image,x2_cen,y2_cen,ref2_r1,ref2_r2))
                    aperture_measurement_ref3.append(aperture(image, x3_cen,y3_cen,radius_ref3))
                    annulus_measurement_ref3.append(annulus(image,x3_cen,y3_cen,ref3_r1,ref3_r2))
                    aperture_measurement_sn.append(aperture(image,sn_x_cen,sn_y_cen,radius_sn,plot=False))
                    annulus_measurement_sn.append(annulus(image,sn_x_cen,sn_y_cen,sn_r1,sn_r2,plot=False))
                    
                    # Get annulus uncertainties
                    annulus_uncertainty_ref1.append(annulus(image,x1_cen,y1_cen,ref1_r1,ref1_r2,plot=False)[1])
                    annulus_uncertainty_ref2.append(annulus(image,x2_cen,y2_cen,ref2_r1,ref2_r2,plot=False)[1])
                    annulus_uncertainty_ref3.append(annulus(image,x3_cen,y3_cen,ref3_r1,ref3_r2,plot=False)[1])
                    annulus_uncertainty_sn.append(annulus(image,sn_x_cen,sn_y_cen,sn_r1,sn_r2,plot=False)[1])
                    
                    # # Draw centroid target lines
                    # fig = plt.figure(figsize=(10,10))
                    # ax = fig.add_subplot(111, projection=w)
                    
                    # # Zoom on ref 1
                    # ax.set_xlim(x1_cen-100,x1_cen+100)
                    # ax.set_ylim(y1_cen-100,y1_cen+100)
                    # ax.axvline(x1_cen,color="green")
                    # ax.axhline(y1_cen,color="green")
                    # ax.imshow(image, vmin=-100, vmax = 1e2, origin = 'lower', cmap = 'inferno')
                    # plt.savefig("plt/centroid.pdf")
                    # plt.show()
                    
    measured_indices = good_indices[measured_indices_in_good_imgs]
\end{verbatim}

\subsection{First Detection, Calculating Magnitudes and Uncertainties}
\begin{verbatim}
    print("{f} images out of {i} remaining after final filter".format(f=len(measured_indices),i=len(images)))
    ref1_intensities = np.abs(np.array(aperture_measurement_ref1) - np.array(annulus_measurement_ref1)[:, 0])
    ref2_intensities = np.abs(np.array(aperture_measurement_ref2) - np.array(annulus_measurement_ref2)[:, 0])
    ref3_intensities = np.abs(np.array(aperture_measurement_ref3) - np.array(annulus_measurement_ref3)[:, 0])
    sn_signal = np.array(aperture_measurement_sn)
    sn_noise = np.array(annulus_measurement_sn)[:, 0]
    sn_intensities = np.abs(sn_signal - sn_noise)
    sn_signal_noise = sn_signal / sn_noise
    
    # Source Uncertainties
    gain = 4
    sn_source_uncertainty = np.sqrt(sn_intensities/gain)
    ref1_src_uncertainty = np.sqrt(ref1_intensities/gain)
    ref2_src_uncertainty = np.sqrt(ref2_intensities/gain)
    ref3_src_uncertainty = np.sqrt(ref3_intensities/gain)
    
    # Calculate total intensity uncertainties
    sn_intensity_uncertainty = np.square(np.sqrt(np.square(sn_source_uncertainty) + np.square(annulus_uncertainty_sn)))
    ref1_intensity_uncertainty = np.square(np.sqrt(np.square(ref1_src_uncertainty) + np.square(annulus_uncertainty_ref1)))
    ref2_intensity_uncertainty = np.square(np.sqrt(np.square(ref2_src_uncertainty) + np.square(annulus_uncertainty_ref2)))
    ref3_intensity_uncertainty = np.square(np.sqrt(np.square(ref3_src_uncertainty) + np.square(annulus_uncertainty_ref3)))
    
    # Find where s/n ratio goes above 3 for first detection threshold
    detected_sn_3 = np.where(abs(sn_signal_noise) > 3)[0][0]
    reversed_indices = np.where(abs(sn_signal_noise[::-1]) > 3)[0]
    last_detected_sn_3 = len(sn_signal_noise) - 1 - reversed_indices[0]
    reversed_indices = np.where(sn_signal_noise[::-1] > 3)[0]
    non_detected_sn_3 = len(sn_signal_noise) - 1 - reversed_indices[0]
    
    first_detection_fn = sorted_names[measured_indices[detected_sn_3]]
    non_detection_fn = sorted_names[measured_indices[non_detected_sn_3]]
    date_detected_header = get_header(first_detection_fn)['DATE-OBS']
    non_detected_header = get_header(non_detection_fn)['DATE-OBS']
    date_detected = datetime.fromisoformat(date_detected_header)
    non_detected = datetime.fromisoformat(non_detected_header)
    sn_month_day = "{m}-{d}".format(m=date_detected.month,d=date_detected.day)
    sn_non_detected_month_day = "{m}-{d}".format(m=non_detected.month,d=non_detected.day)
    
    print("First Detection: {d}".format(d=date_detected))
    print("Last Detection: {d}".format(d=non_detected))
    # Detection uncertainty
    detection_uncertainty = np.sqrt(np.square(sn_source_uncertainty/sn_signal) + np.square(annulus_uncertainty_sn/sn_noise))*(sn_signal/sn_noise)
    print("detected uncertainty: {u}".format(u=detection_uncertainty[detected_sn_3]))
    print("last detection uncertainty {u}".format(u=detection_uncertainty[last_detected_sn_3]))
    
    # Calculate magnitudes
    ref3_mag = 15.61 # Uncertainty 0.02
    ref2_mag = -2.5*np.log10(ref2_intensities/ref3_intensities) + ref3_mag
    ref1_mag = -2.5*np.log10(ref1_intensities/ref3_intensities) + ref3_mag
    sn_mag = 2.5*np.log10(ref3_intensities/sn_intensities) + ref3_mag
    
    # Calculate residuals of reference stars
    ref1_expected = 16.32 # uncertainty 0.07
    ref2_expected = 17.16 # uncertainty 0.08
    
    # Using ref 3 as the reference star
    ref1_residuals = abs(ref1_mag - ref1_expected)
    ref1_median_residual = np.median(ref1_residuals)
    ref2_residuals = abs(ref2_mag - ref2_expected)
    ref2_median_residual = np.median(ref2_residuals)
    measurements = np.linspace(0,len(ref1_residuals),len(ref1_residuals))
    
    # Check residuals of REF 1 and REF 2
    bounds = 1.3
    fig, ax = plt.subplots(2, 1, figsize=(6, 8), sharex=True)
    ax[0].set_xlabel("Measurement Number")
    ax[0].scatter(measurements, ref1_residuals, color="red", s=2, label="REF 1 Median Residual: {r}".format(r=round(ref1_median_residual, 2)))
    ax[0].set_title("Residuals of Reference Stars")
    ax[0].set_ylabel("Residuals")
    ax[0].set_ylim(-bounds, bounds)
    ax[0].grid(True)
    ax[1].scatter(measurements, ref2_residuals, color="teal", s=2, label="REF 2 Median Residual: {r}".format(r=round(ref2_median_residual, 2)))
    ax[1].set_ylabel("Residuals")
    ax[1].set_ylim(-bounds, bounds)
    ax[1].grid(True)
    plt.tight_layout()
    plt.subplots_adjust(hspace=0)
    ax[0].legend(loc='best', fontsize=14)
    ax[1].legend(loc='best', fontsize=14)
    plt.savefig("plt/residuals.pdf")
    plt.show()
    
    # Propogate reference star uncertainties to get SN magnitude uncertainties
    ref3_mag_uncertainty = 0.02
    sigma_ref1 = (2.5/np.log(10)) * (sn_intensity_uncertainty/sn_intensities)
    sigma_ref2 = (2.5/np.log(10)) * (ref2_intensity_uncertainty/ref2_intensities)
    sn_magnitude_uncertainty = np.sqrt(sigma_ref1**2 + sigma_ref2**2 + ref3_mag_uncertainty**2)
    
    plt.title("Supernova Magnitudes")
    plt.xlabel("Image Number")
    plt.ylabel("B-Band Magnitude")
    plt.ylim(17.5, 24)
    plt.scatter(measured_indices, sn_mag, color="blue", label="Supernova Magnitude", s=2)
    plt.errorbar(measured_indices, sn_mag, yerr=sn_magnitude_uncertainty, fmt='none', color='g', label='Uncertainty', lw=0.5, alpha=0.6, ls='-')
    plt.axvline(measured_indices[detected_sn_3], color="red", ls="--", label="Epoch of First Detection: {i}".format(i=date_detected))
    plt.axvline(measured_indices[non_detected_sn_3], color="green", ls="--", label="Epoch of Last Detection: {i}".format(i=non_detected))
    plt.gca().invert_yaxis()
    plt.grid(".-")
    plt.legend()
    plt.savefig("plt/raw_mag.pdf")
    plt.show()
\end{verbatim}

\subsection{Bin Images by Date}
\begin{verbatim}
    # Get the dates of all measured images
    month_days = []
    full_dates = []
    for index in measured_indices:
        # Get date
        current_name = sorted_names[index]
        current_header = get_header(current_name)
        current_date = datetime.fromisoformat(current_header['DATE-OBS'])
        month_days.append("{month}-{date}".format(month=current_date.month, date=current_date.day))
        full_dates.append(current_date)
    
    # Bin the images by day
    days = []
    measurements = []
    uncertainties = []
    full_names = []
    intensities = []
    for i in range(len(month_days)):
        day = month_days[i]
        full_name = full_dates[i]
        measurement = sn_mag[i]
        intensity_current = sn_intensities[i]
        uncertainty = sn_magnitude_uncertainty[i]
        if day not in days:
            days.append(day)
            measurements.append([measurement])
            uncertainties.append([uncertainty])
            full_names.append([full_name])
            intensities.append([intensity_current])
        else:
            index = days.index(day)
            measurements[index].append(measurement)
            uncertainties[index].append(uncertainty)
            full_names[index].append(full_name)
            intensities[index].append(intensity_current)
    
    # Plot the binned magnitudes with uncertainties
    fig = plt.figure(figsize=(12,8))
    plt.title("Supernova Magnitudes")
    plt.xlabel("Days")
    plt.ylabel("B-Band Magnitude")
    plt.ylim(17.5, 24.5)
    plt.xticks(rotation=45)
    plt.gca().invert_yaxis()
    for i in range(len(days)):
        day = days[i]
        measurement_list = measurements[i]
        uncertainty_list = uncertainties[i]
        measurements_in_day = len(measurement_list) * [day]
        plt.errorbar(measurements_in_day, measurement_list, yerr=uncertainty_list, fmt='o', markersize=3, capsize=3, alpha=0.3, lw=0.5, color='g', ls='--', ecolor='g')
        plt.scatter(measurements_in_day, measurement_list, color='blue', s=5)
    detection = plt.axvline(sn_month_day, color="red", ls="--", label="First Detection: {i}".format(i=date_detected))
    mag_label = lines.Line2D([], [], color='blue', marker='o', markersize=3, ls='None', label="Magnitude Measurement")
    uncertainty_label = lines.Line2D([], [], color='darkblue', marker='o', markersize=3, ls='None', label="Uncertainty")
    plt.legend(loc="lower right", fontsize=14, handles=[mag_label, uncertainty_label, detection])
    plt.grid(".-")
    fig.savefig("plt/binned_uncertainty.pdf")
    plt.show()
\end{verbatim}

\subsection{Weighted Mean Polynomial Fitting, Peak Magnitude Finding}
\begin{verbatim}
    # Calculate weighted mean of each day based on uncertainties, also calculate mean uncertainty
    weighted_means = []
    mean_uncertainties = []
    for i in range(len(days)):
        measurement_list = measurements[i]
        uncertainty_list = uncertainties[i]
        weights = 1 / np.array(uncertainty_list)
        weighted_mean = np.sum(measurement_list * weights) / np.sum(weights)
        mean_uncertainty = 1/np.sum(1/np.array(uncertainty_list)**2)
        weighted_means.append(weighted_mean)
        mean_uncertainties.append(mean_uncertainty)
    
    fig = plt.figure(figsize=(12,8))
    plt.title("Weighted Mean of Supernova Magnitudes")
    plt.xlabel("Days")
    plt.ylabel("B-Band Magnitude")
    plt.ylim(17.5, 24.5)
    plt.xticks(rotation=45)
    plt.gca().invert_yaxis()
    plt.scatter(days, weighted_means, color='blue', label='Weighted Mean', s=40)
    plt.errorbar(days, weighted_means, yerr=mean_uncertainties, fmt='none', color='green', label='Mean Uncertainty', lw=0.5, alpha=0.3, ls='--')
    plt.axvline(sn_month_day, color="red", ls="--", label="First Detection: {i}".format(i=date_detected))
    plt.axvline(sn_non_detected_month_day, color="brown", ls="--", label="Non-Detection: {i}".format(i=non_detected))
    mag_label = lines.Line2D([], [], color='blue', marker='o', markersize=3, ls='None', label="Weighted Mean")
    uncertainty_label = lines.Line2D([], [], color='green', marker='o', markersize=3, ls='None', label="Mean Uncertainty")
    detection = plt.axvline(sn_month_day, color="red", ls="--", label="First Detection: {i}".format(i=date_detected))
    non_detection = plt.axvline(sn_non_detected_month_day, color="brown", ls="--", label="Non-Detection: {i}".format(i=non_detected))
    plt.legend(loc="lower right", fontsize=14, handles=[mag_label, uncertainty_label, detection, non_detection])
    plt.grid(".-")
    fig.savefig("plt/weighted_mean.pdf")
    plt.show()
    
    # Fit a polynomial to the weighted mean
    poly2 = np.polyfit(np.arange(len(days)), weighted_means, 2)
    poly3 = np.polyfit(np.arange(len(days)), weighted_means, 3)
    poly4 = np.polyfit(np.arange(len(days)), weighted_means, 4)
    
    # Get uncertainty of the peak magnitude
    peak_magnitude_uncertainty2 = np.std(abs(weighted_means - np.polyval(poly2, np.arange(len(days)))))
    peak_magnitude_uncertainty3 = np.std(abs(weighted_means - np.polyval(poly3, np.arange(len(days)))))
    peak_magnitude_ucnertainty4 = np.std(abs(weighted_means - np.polyval(poly4, np.arange(len(days)))))
    
    # Evaluate magnitudes from polynomials and find the peak magnitude
    magnitudes4 = np.polyval(poly4, np.arange(len(days)))
    magnitudes3 = np.polyval(poly3, np.arange(len(days)))
    magnitudes2 = np.polyval(poly2, np.arange(len(days)))
    
    peak_day_index3 = np.argmin(magnitudes3)
    peak_day_index2 = np.argmin(magnitudes2)
    peak_day_index4 = np.argmin(magnitudes4)
    
    peak_time_measurement3 = np.argmax(np.array(measurements[peak_day_index3]))
    peak_time_measurement2 = np.argmax(np.array(measurements[peak_day_index2]))
    peak_time_measurement4 = np.argmax(np.array(measurements[peak_day_index4]))
    
    peak_time_name4 = full_names[peak_day_index4][peak_time_measurement4]
    peak_time_name3 = full_names[peak_day_index3][peak_time_measurement3]
    peak_time_name2 = full_names[peak_day_index2][peak_time_measurement2]
    
    peak_magnitude3 = magnitudes3[peak_day_index3]
    peak_magnitude2 = magnitudes2[peak_day_index2]
    peak_magnitude4 = magnitudes4[peak_day_index4]
    
    date_detected_index = days.index(sn_month_day)
    date_detected_unbinned = month_days.index(sn_month_day)
    
    
    print("Peak Magnitude (Quadratic): {m}".format(m=peak_magnitude2))
    print("Peak Magnitude (Cubic): {m}".format(m=peak_magnitude3))
    print("Peak Magnitude (Quartic): {m}".format(m=peak_magnitude4))
    print("--------------------------------")
    print("Peak Magnitude Uncertainty (Linear): {m}".format(m=peak_magnitude_uncertainty2))
    print("Peak Magnitude Uncertainty (Cubic): {m}".format(m=peak_magnitude_uncertainty3))
    print("Peak Magnitude Uncertainty (Quartic): {m}".format(m=peak_magnitude_ucnertainty4))
\end{verbatim}

\subsection{Plot Full Beautiful Curve}
\begin{verbatim}
days_since_detection = [i - date_detected_index - 1 for i in range(len(days))]
days_since_detection_np = np.array(days_since_detection)

# Plot weighted mean along with binned data, and fit polynomials to the weighted mean.
fig, ax1 = plt.subplots(figsize=(12,8))
ax1.set_title("Supernova Magnitude", pad=5)
ax1.set_xlabel("Date")
ax1.set_ylabel("B-Band Magnitude")
ax1.set_ylim(17.5, 24.5)
ax1.set_xticks(np.arange(len(days)))
ax1.set_xticklabels(days, rotation=45)
ax1.invert_yaxis()
for i in range(len(days)):
    day = days[i]
    measurement_list = measurements[i]
    uncertainty_list = uncertainties[i]
    measurements_in_day = len(measurement_list) * [day]
    ax1.errorbar(measurements_in_day, measurement_list, yerr=uncertainty_list, fmt='o', markersize=3, capsize=3, alpha=0.3, lw=0.5, color='blue', ls='--', ecolor='darkblue')
    ax1.scatter(measurements_in_day, measurement_list, color='blue', s=5)
ax1.scatter(days, weighted_means, color='red', marker='x', label='Weighted Mean', s=50)
ax2 = ax1.twiny()
ax2.set_xlabel("Days Since Epoch of First Detection", fontsize=12)
ax2.set_xticks(days_since_detection_np)
ax2.set_xticklabels(days_since_detection, rotation=45)
ax2.xaxis.set_ticks_position('top')
ax1.plot(days, magnitudes4, color='orange', label='4th Order Polynomial Fit', ls='--')
non_detection = ax1.axvline(sn_non_detected_month_day, color="green", ls="--", label="Epoch of Non-Detection: {i}".format(i=non_detected))
detection = ax1.axvline(sn_month_day, color="red", ls="--", label="Epoch of First Detection: {i}".format(i=date_detected))
peak_brightness_quartic = ax1.axvline(days[peak_day_index4], color="blue", ls="--", label="Epoch of Peak Brightness: {i}".format(i=peak_time_name4))
mag_label = lines.Line2D([], [], color='blue', marker='o', markersize=3, ls='None', label="Magnitude Measurement")
uncertainty_label = lines.Line2D([], [], color='darkblue', marker='o', markersize=2, ls='None', label="Uncertainty")
weighted_mean_label = lines.Line2D([], [], color='red', marker='x', markersize=9, ls='None', label="Weighted Mean")
polynomial_label = lines.Line2D([], [], color='orange', ls='--', label="4th Order Polynomial Fit")
plt.legend(loc="lower right", fontsize=14, handles=[mag_label, uncertainty_label, weighted_mean_label, polynomial_label, detection, non_detection, peak_brightness_quartic])
plt.grid(".-")
plt.savefig("plt/weighted_mean_poly.pdf")
plt.show()
\end{verbatim}

\subsection{Normalized Intensity Power Law Fitting}
\begin{verbatim}
    normalized_intensity = 10**(-(np.array(sn_mag) - peak_magnitude4)/2.5)
    # Plot normalized light curve
    t = np.arange(0, len(normalized_intensity[date_detected_unbinned:]),1)
    t_days = np.arange(0,len(days))
    plt.xticks(t_days)
    plt.title("Normalized Intensity Supernova Curve")
    plt.scatter(t, normalized_intensity[date_detected_unbinned:], color="red", label="Normalized Intensity", s=5)
    plt.xlabel("Days Since Epoch of First Detection")
    plt.ylim(0,1.2)
    plt.ylabel("Normalized Intensity")
    plt.grid(".-")
    plt.legend()
    plt.show()
    
    def func(t, C, t1):
        return C*(t-t1)**2
    
    # data to fit
    data = normalized_intensity[:date_detected_unbinned+20]
    t_data = np.arange(np.negative(len(normalized_intensity[:date_detected_unbinned])), len(normalized_intensity[date_detected_unbinned:date_detected_unbinned+20]), 1)
    
    # Initial guess for the parameters
    C_initial = 0.004279039146597766
    t1_initial = -1
    p0 = [C_initial, t1_initial]
    params, params_covariance = curve_fit(func, t_data, data, p0)
    
    # Results
    print("C =", params[0])
    print("t1 =", params[1])
    
    # Calculate the residual and the reduced chi-square
    residuals = data - func(t_data, params[0], params[1])
    chi_squared = np.sum((residuals)**2)
    reduced_chi_squared = chi_squared / (len(t_data) - len(p0))
    
    print("Reduced Chi-squared: ", reduced_chi_squared)
    print("Covariance Matrix: ", params_covariance)
    print("Uncertainties in parameters: ", np.sqrt(np.diag(params_covariance)))
    
    # Plot the fitted curve
    plt.title("Normalized Intensity Supernova Curve)")
    plt.xticks(t_days)
    plt.scatter(t, normalized_intensity[date_detected_unbinned:], color="red", label="Normalized Intensity", s=5)
    plt.plot(t_data, func(t_data, params[0], params[1]), label='Chi-Squared = {c}'.format(c=round(reduced_chi_squared,4)))
    plt.xlabel("Days")
    plt.ylim(0,0.4)
    plt.xlim(0,19)
    plt.xticks(rotation=45) 
    plt.ylabel("Normalized Intensity")
    plt.grid(".-")
    plt.legend()
    plt.savefig("plt/power_law.pdf")
    plt.show()
\end{verbatim}

\end{appendices}

\end{document}